\section{Products}
\label{sec:products}

All three products share a single protocol layer (the underwriter, per-market order book, settlement pipeline, and risk engine), extensible to new products without modifying core contracts.

\subsection{Borrow}

A holder of any supported token borrows against it with configurable rate models (fixed-rate or variable-rate referencing an external index) for a defined term. Collateral stays on its home chain; yield-bearing collateral (e.g.\ wstETH) continues accruing native yield throughout the loan. Loan-to-value ratio, rate, and term are determined per loan through operator competition (Section~\ref{sec:architecture}). Origination, interest accrual, and liquidation are each governed by zero-knowledge programs, ensuring that no single trusted party controls the loan's lifecycle.

\subsection{Prime Trade}

A borrower draws a credit line in the market's supported principal asset and routes it to a per-loan margin agent on HyperEVM, gaining leveraged perpetual exposure without selling the underlying collateral (Section~\ref{sec:crosschain}). In the initial markets, USDC is the canonical principal asset because of its deep DeFi liquidity. USDC is therefore the reference implementation, not a permanent protocol constraint: other stablecoins may be supported by separate markets if the relevant liquidity, risk, and compliance requirements are satisfied.

The margin agent restricts capital to approved operations (trading only, no withdrawals), so the lender's risk depends on programmatic constraints, not borrower discretion. Venue leverage multiplies the borrowing LTV, controlling large notional with modest collateral (Section~\ref{sec:crosschain}).

\subsection{Earn}

Liquidity providers grant an allowance to an underwriter and receive ERC-20 shares representing their proportional claim on deployed capital and returns. Proof-gated settlement deploys underwriter capital into loans originated through the protocol's order book, either funding loan principal directly or providing hedging buffers when loans are sourced through external venues (Section~\ref{sec:capital}). The underwriter's strategy is bounded by a committed policy hash; compliance is enforced at proof time (Section~\ref{sec:lp}). Underwriter policy selects the rate model: under fixed-rate terms the coupon is locked at origination; under variable-rate terms an external index is passed through to borrowers. On direct-funded loans, yield derives from the full coupon; on liquidity-forwarding loans, from the hedging-buffer surplus. LP funds leave their wallets only at the instant of proof-gated deployment into a specific position.
